{\rtf1\ansi\ansicpg1252\cocoartf2870
\cocoatextscaling0\cocoaplatform0{\fonttbl\f0\fswiss\fcharset0 Helvetica;}
{\colortbl;\red255\green255\blue255;}
{\*\expandedcolortbl;;}
\paperw11900\paperh16840\margl1440\margr1440\vieww11520\viewh8400\viewkind0
\pard\tx720\tx1440\tx2160\tx2880\tx3600\tx4320\tx5040\tx5760\tx6480\tx7200\tx7920\tx8640\pardirnatural\partightenfactor0

\f0\fs24 \cf0 \\documentclass[11pt,a4paper]\{article\}\
\\usepackage[utf8]\{inputenc\}\
\\usepackage\{amsmath, amssymb\}\
\\usepackage\{geometry\}\
\\geometry\{a4paper, margin=1in\}\
\\usepackage\{hyperref\}\
\\usepackage\{graphicx\}\
\\documentclass\{article\}\
\\usepackage\{booktabs\}\
\\usepackage\{tabularx\}\
\\usepackage[margin=1in]\{geometry\}\
% !TEX program = pdflatex\
\\documentclass[11pt,a4paper]\{article\}\
\\usepackage[utf8]\{inputenc\}\
\\usepackage[T1]\{fontenc\}\
\\usepackage\{lmodern\}\
\\usepackage\{geometry\}\
\\usepackage\{hyperref\}\
\\usepackage\{enumitem\}\
\\geometry\{margin=2.5cm\}\
\
\
\\begin\{document\}\
\\section*\{Learning Objectives of the Thesis\}\
This thesis provides an opportunity for me to explore methodologies for transforming black-box AI models into white-box systems. A primary objective is to learn how to mathematically model AI systems resulting in formal representations of these systems. This work aims to enhance the understanding of how RAG systems react to various system prompts. Since system prompts are plain english text, the research questions requires a treatment of natural language. A subsequent goal is to refine skills in selecting appropriate experimental designs, conducting experiments based on these designs, and drawing conclusions from the results as well as writing everything up in a concise way.\
\
\\section*\{Structure of the Thesis (Strucutre Inspired by Moritz's Thesis)\}\
\\begin\{description\}[font=\\bfseries, style=nextline, leftmargin=1.5em]\
    \\item[1. Introduction (approx. 2 pages)] \
    Provides the background, motivation, research objectives, and an overview of the thesis structure.\
    \
    \\item[2. Preliminaries (approx. 6 pages)] \
    Offers a comprehensive overview of RAG systems detailed across multiple subsections.\
    \
    \\item[3. Literature Review (approx. 3 pages)] \
    Summarizes the existing research relevant to the study.\
    \
    \\item[4. Modeling Approach and Experimental Setup (approx. 5 pages)] \
    Details the mathematical modeling of the language and RAG systems, as well the experimental design.\
    \
    \\item[5. Results and Discussion (approx. 5 pages)] \
    Presents the empirical findings from the experiments and discusses their implications.\
    \
    \\item[6. Conclusion (approx. 1 page)] \
    Concludes the thesis by summarizing findings and suggesting potential directions for future work.\
    \
    \\item[Appendix] \
    Contains supplementary materials.\
\\end\{description\}\
\\section*\{Elevator Pitch for the Impatient\}\
\\textbf\{Central research question:\} To what extent do systematic changes in the system prompt alter the behavior and outcomes of a RAG system when the user input, knowledge base, retrieval mechanism, and underlying language model are held constant?\\\\\
\\newline\
\\textbf\{Objective:\} study the effects of different system prompts on system behavior and system outcomes within a simple RAG system. \\\\\
\\newline\
\\textbf\{Background:\} RAG systems combine information retrieval with language generation. A user query is first used to retrieve relevant information from an external knowledge base (Vector Database). The retrieved information is subsequently provided to a language model together with instructions for the generation of the final response. The system prompt provides an important interface between the intended behavior of the system and its observable output.\\\\\
\\newline\
\\textbf\{Related literature (Summaries by ChatGPT and refined by me):\}\\\\\
\\newline\
\\noindent\\textit\{Cheng and Mastropaolo (2026)\}: study how different system prompts affect instruction-tuned language models for code generation, varying prompt specificity, prompting strategy, model, programming language, and decoding settings. Their work directly motivates the controlled manipulation of system prompts in this thesis, while the present experiment extends the analysis to a RAG setting and focuses explicitly on system behaviour and system outcomes.\
\
\\medskip\
\\noindent\\textit\{Mu et al. (2025)\}:\
investigate the robustness of system prompts as a mechanism for controlling LLM behaviour, particularly when system instructions interact with conflicting or adversarial user inputs. Their work supports modelling the system prompt as a behavioural control variable and motivates measuring whether an system actually follows the intended instructions.\
\
\\medskip\
\\noindent\\textit\{Lu et al. (2024)\}: study prompt sensitivity across different models and tasks, examining how changes in prompt formulation affect model behaviour and output. This is relevant to the present research because it motivates distinguishing changes in the intended instruction from changes caused by the linguistic formulation of that instruction.\
\
\\medskip\
\\noindent\\textit\{Zhuo et al. (2024)\}: propose methods for measuring and understanding prompt sensitivity across LLMs and tasks, investigating why seemingly small prompt changes can produce different model behaviour. Their work provides methodological support for treating prompt effects as measurable experimental phenomena rather than anecdotal observations.\
\
\\medskip\
\\noindent\\textit\{Salinas and Morstatter (2024)\}: The authors investigate the sensitivity of LLM decisions to small perturbations in prompts, demonstrating that minor changes in prompt construction can sometimes produce substantial changes in model outputs. This motivates the use of controlled, minimal prompt interventions in my thesis.\
\
\\medskip\
\\noindent\\textit\{Zheng et al. (2024)\}: investigate the effects of assigning different personas or roles through system prompts and examine whether these changes improve model performance. Their findings are relevant to the distinction between changes in intended system behaviour and actual improvements in objective system performance.\
\
\\medskip\
\\noindent\\textit\{Ying et al. (2024)\}: examine how LLMs respond to conflicting prompts and information, including situations relevant to retrieval-augmented generation. Their work is particularly relevant to the thesis' interest in how system instructions influence an system when retrieved evidence is incomplete, conflicting, or inconsistent with other information.\
\
\\medskip\
\\noindent\\textit\{Ruiz and Cardona-Betancourt (2026)\}:  experimentally investigate how LLMs respond to explicit behavioural rules provided through prompts and how these rules interact with changing task contingencies. This supports the proposed view of system prompts as specifications of behavioural policies rather than only textual context.\
\
\\medskip\
\\noindent\\textit\{Hua et al. (2025)\}: examine prompt sensitivity and show that some apparent differences between prompts can arise from the way model outputs are evaluated rather than from meaningful behavioural differences. This is important for the present research because it motivates using evaluation metrics that distinguish genuine changes in RAG behaviour and outcomes from superficial output variation.\
\
\\medskip\
\\noindent\\textit\{Qin et al. (2026)\}: investigate prompt sensitivity through interactions between prompts and model inputs, moving beyond simple comparisons of final outputs to analyze how prompt changes affect model behaviour. Their approach is closely related to the present research's goal of developing a more structured account of how system prompts influence an AI system.\\\\\
\\newline\
\\textbf\{Methodological challenge:\} isolate the effect of the system prompt from other sources of variation. The model should contain only those variables and transformations that are necessary to answer the research question. The proposed architecture deliberately excludes unnecessary system memory, multi-system interaction, query planning, and complex tool-use mechanisms etc. \\\\\
\\newline\
\\textbf\{Modeling approach:\} the model treats the RAG system as a composition of a small number of functions. This provides a controlled experimental environment in which the system prompt can be treated as the independent variable while other relevant components remain fixed.\\\\\
\\newline\
experimental abstraction of RAG:\
\\[\
\\boxed\{\
F(q;S,D) = G(q,R(q;D);S)\
\}\
\\]\
with\
\\[\
q\\in\\mathcal Q, \\qquad S\\in\\mathcal S, \\qquad D\\in\\mathcal D, \\qquad C=R(q;D), \\qquad y\\in\\mathcal Y.\
\\]\
\\newline\
The system prompt is parameterized as\
\\[\
S=S_\\theta,\
\\]\
where $\\theta$ represents the experimentally controlled instruction, constraint, and output requirements.\\\\\
\\newline\
\\textbf\{Experimental approach:\} The experimental procedure is:\
\\[\
\\boxed\{\
\\begin\{array\}\{c\}\
\\text\{Construct controlled knowledge base\}\\\\\
\\downarrow\\\\\
\\text\{Construct query set\}\\\\\
\\downarrow\\\\\
\\text\{Establish baseline \}S_0\\\\\
\\downarrow\\\\\
\\text\{Apply controlled prompt intervention \}S_0\\rightarrow S_j\\\\\
\\downarrow\\\\\
\\text\{Hold retrieval and model configuration fixed\}\\\\\
\\downarrow\\\\\
\\text\{Measure behavior and outcomes\}\\\\\
\\downarrow\\\\\
\\text\{Compare \}\\Delta M_j\
\\end\{array\}\
\}\
\\]\
\\newline\
The fundamental research relationship is \
\\[\
\\boxed\{\
S \\rightarrow \\text\{system Behavior\} \\rightarrow \\text\{System Outcome\}.\
\}\
\\]\
\\newline\
\\newpage\
\\section*\{Formal Modeling of the RAG System\}\
Let\
\\[\
q \\in \\mathcal\{Q\}\
\\]\
denote a user input, where $\\mathcal\{Q\}$ is the space of possible user queries.\\\\\
\\newline\
Let\
\\[\
D=\\\{d_1,d_2,\\ldots,d_N\\\}\
\\]\
denote the knowledge base, where each $d_i$ is a document and $N$ is the total number of documents.\\\\\
\\newline\
Let\
\\[\
S \\in \\mathcal\{S\}\
\\]\
denote a system prompt, where $\\mathcal\{S\}$ is the space of possible system prompts.\\\\\
\\newline\
Let\
\\[\
C \\in \\mathcal\{C\}\
\\]\
denote the retrieved context supplied to the language model.\\\\\
\\newline\
Finally, let\
\\[\
y \\in \\mathcal\{Y\}\
\\]\
denote the generated system output, where $\\mathcal\{Y\}$ is the space of possible responses.\\\\\
\\newline\
The simple RAG system can be represented by the following sequence:\
\\[\
q \\longrightarrow C \\longrightarrow y.\
\\]\
\\newline\
The system prompt influences the generation stage:\
\\[\
(q,C,S) \\longrightarrow y.\
\\]\
\\newline\
The complete system can be represented as\
\\[\
F(q;S,D)=G(q,R(q;D);S),\
\\]\
where $R$ is the retrieval function and $G$ is the generation function.\
\\newline\
\\subsection*\{Retrieval Function\}\
\\newline\
The Retriever is defined as a function\
\\[\
R:\\mathcal\{Q\}\\times\\mathcal\{D\}\\rightarrow\\mathcal\{C\},\
\\]\
where $\\mathcal\{D\}$ denotes the space of knowledge bases and $\\mathcal\{C\}$ denotes the space of retrieved contexts.\\\\\
\\newline\
For a query $q$ and knowledge base $D$,\
\\[\
C=R(q;D).\
\\]\
\\newline\
The retrieved context may consist of the top $k$ documents according to a relevance function $r$:\
\\[\
C=R_k(q;D) = \\operatorname\{TopK\}_\{d\\in D\} r(q,d).\
\\]\
\\newline\
The precise implementation of $r$ is not the primary object of this exercise. It is therefore fixed throughout the experiment. This allows changes in system outcomes to be attributed primarily to changes in the system prompt.\\\\\
\
\\subsection*\{Generation Function\}\
\
The Generator is defined as\
\\[\
G:\\mathcal\{Q\}\\times\\mathcal\{C\}\\times\\mathcal\{S\}\\rightarrow\\mathcal\{Y\}.\
\\]\
\\newline\
Given a user query $q$, retrieved context $C$, and system prompt $S$, the generator produces\
\\[\
y=G(q,C;S).\
\\]\
\\newline\
The generator includes the underlying language model and the prompt construction required to provide the model with the user query, retrieved context, and system instructions.\\\\\
\\newline\
The complete RAG system is therefore\
\\[\
\\boxed\{\
F(q;S,D) = G\\left(q,R(q;D);S\\right)\
\}\
\\]\
\\newline\
This representation is as minimal as possible:\
\\[\
\\boxed\{\
\\text\{Input\} \\rightarrow \\text\{Retrieval\} \\rightarrow \\text\{Generation\} \\rightarrow \\text\{Output\}\
\}\
\\]\
with the system prompt serving as a controlled parameter of generation.\\\\\
\
\\subsection*\{Controlled Variables\}\
\
The exercise's goal is to estimate the effect of the system prompt $S$ on the output. The following variables should remain constant across the different experiments:\
\\[\
q,\\quad D,\\quad R,\\quad G.\
\\]\
\\newline\
The experimental treatment is the system prompt:\
\\[\
S=S_j,\
\\]\
where $j$ indexes a particular prompt condition.\\\\\
\\newline\
For a fixed query and knowledge base,\
\\[\
y_j = F(q;S_j,D).\
\\]\
\\newline\
The comparison between $y_i$ and $y_j$  represents the observed effect of changing the system prompt while keeping the remaining system configuration constant.\
\
\\section*\{Modeling the System Prompt\}\
I define the system prompt as a parameter vector\
\\[\
\\theta=(\\theta_I,\\theta_C,\\theta_O),\
\\]\
where:\
\\[\
\\theta_I \\quad \\text\{represents the selected instructions,\}\
\\]\
\\[\
\\theta_C \\quad \\text\{represents the selected constraints, and\}\
\\]\
\\[\
\\theta_O \\quad \\text\{represents the selected output requirements.\}\
\\]\
\\newline\
The system prompt can then be represented as\
\\[\
S_\\theta.\
\\]\
\\newline\
The whole RAG system can be written as\
\\[\
\\boxed\{\
y_\\theta = G(q,R(q;D);S_\\theta)\
\}\
\\]\
\\newline\
and the experimental objective becomes the comparison of\
\\[\
y_\{\\theta_1\},y_\{\\theta_2\},\\ldots,y_\{\\theta_m\}\
\\]\
\\newline\
keeping all other variables fixed.\
\\subsubsection*\{Linguistic Realization\}\
\
Since the thesis concerns system prompts, the linguistic form of the instructions should also be controlled.\\\\\
\\newline\
Let\
\\[\
L(S_\\theta)\
\\]\
denote the the particular wording used to express an instruction of a system prompt.\\\\\
\\newline\
Two system prompts may have approximately equivalent functional content while differing linguistically:\
\\[\
L(S_\{\\theta_1\})\\neq L(S_\{\\theta_0\}),\
\\]\
while being semantically equivalent\
\\[\
\\operatorname\{Sem\}(S_\{\\theta_1\})\\approx\\operatorname\{Sem\}(S_\{\\theta_0\}).\
\\]\
\\newline\
For the primary experiment, however, linguistic variation should not be introduced simultaneously with changes in semantics. The primary experimental comparison will be using minimal prompt perturbations. This creates a controlled intervention:\
\\[\
S_\{\\theta_0\} \\rightarrow (S_\{\\theta_1\}=S_\{\\theta_0\}+\\Delta S)\
\\]\
\\newline\
The effect of the intervention is then measured as\
\\[\
\\Delta Y = M(y_\{\\theta_1\})-M(y_\{\\theta_0\}),\
\\]\
where $M$ is an evaluation metric.\\\\\
\\subsection*\{Variation in System Prompts\}\
\
A small number of different prompt conditions should be sufficient for the scope of this exercise.\
\\begin\{align*\}\
S_\{\\theta_0\} &= \\text\{Canonical. Plain imperative prose. The reference form.\}, \\\\\
S_\{\\theta_1\} &= \\text\{Lexical. Synonym substitution, register held fixed. No syntactic change.\}, \\\\\
S_\{\\theta_2\} &= \\text\{Syntactic. Clause reordering, active\uc0\u8596 passive, imperative\u8596 declarative. Lexicon held as close as possible\}, \\\\\
S_\{\\theta_3\} &= \\text\{Format. Prose \uc0\u8594  bulleted list, punctuation, casing, delimiters. Propositional content byte-identical where possible. This is the Sclar et al. axis \'97 pure spurious features.\}, \\\\\
S_\{\\theta_4\} &= \\text\{Deontic. Modal force: must / should / please try to. This is a deliberate boundary probe. It is arguably semantic, not linguistic. Include it precisely because it tests where the L/Sem line actually falls \'97 and say so in the thesis. If L\uc0\u8324  produces large effects while L\u8321 \'96L\u8323  don't, you've located the boundary empirically rather than stipulating it\}.\
\\end\{align*\}\
The working hypothesis is that these perturbation do not change the sematic meaning of the system prompt\
\\newpage\
\\section*\{Experimental Design\}\
The experiment estimates the effect of the system prompt on observable system behavior and system outcomes.\\\\\
\\newline\
For each query $q_i$ and prompt condition $S_j$, define\
\\[\
y_\{ij\} = F(q_i;S_j,D).\
\\]\
\\newline\
The experimental dataset consists of observations\
\\[\
\\mathcal\{E\} = \\\{(q_i,S_j,C_i,y_\{ij\})\\\},\
\\]\
where\
\\[\
C_i=R(q_i;D).\
\\]\
\\newline\
The same retrieved context should be used across prompt conditions since this is particularly important for isolating the effect of the system prompt.\\\\\
\\subsubsection*\{Step 1: Knowledge Base\}\
A controlled knowledge base $D$ should be constructed containing documents that provide known answers to the experimental queries.\\\\\
\\newline\
The corpus should be comprised of sufficient variation:\
\\begin\{enumerate\}\
    \\item directly answerable questions\
    \\item questions requiring information from multiple documents\
    \\item questions for which the evidence is incomplete\
    \\item questions for which retrieved documents contain conflicting information\
    \\item questions for which irrelevant but superficially similar documents are retrieved\
\\end\{enumerate\}\
\\newline\
The true relationship between questions and supporting documents should be known ex-ante. \
\\subsubsection*\{Step 2: Query Set\}\
Let\
\\[\
Q=\\\{q_1,q_2,\\ldots,q_M\\\}\
\\]\
denote the experimental query set.\\\\\
\\newline\
Each query should be associated with a known answer or evaluation criterion.\\\\\
\\newline\
For example,\
\\[\
Q = Q_\{\\text\{direct\}\} \\cup Q_\{\\text\{multi-hop\}\} \\cup Q_\{\\text\{conflict\}\} \\cup Q_\{\\text\{insufficient\}\}.\
\\]\
This allows the effect of a prompt to be evaluated under different information conditions.\\\\\
\\subsubsection*\{Step 3: Baseline\}\
A baseline system is defined as\
\\[\
F_0(q)=G(q,R(q;D);S_0),\
\\]\
where $S_0$ is the baseline system prompt.\\\\\
\\newline\
Every query is first evaluated using $S_0$. The resulting baseline establishes the reference values for the subsequent experimental comparisons.\\\\\
\\newline\
For each query $q_i$, record:\
\\[\
C_i=R(q_i;D)\
\\]\
and\
\\[\
y_\{i0\}=F_0(q_i).\
\\]\
\\newline\
The baseline should also record the retrieval results, model configuration, decoding parameters, and other implementation details.\\\\\
\\subsubsection*\{Step 4: Prompt Interventions\}\
For each experimental prompt $S_j$, repeat the experiment using the same query set:\
\\[\
y_\{ij\} = G(q_i,R(q_i;D);S_j).\
\\]\
\\newline\
The only intentional perturbation is\
\\[\
S_0\\rightarrow S_j.\
\\]\
\\newline\
This creates a paired experimental design because the same query is evaluated under multiple prompt conditions.\\\\\
\\subsubsection*\{Step 5: Repeat Runs\}\
Language-model generation is stochastic. Therefore, each condition should be repeated $n$ times.\\\\\
\\newline\
Let\
\\[\
y_\{ijr\}\
\\]\
denote the output for query $q_i$, prompt $S_j$, and repetition $r$.\\\\\
\\newline\
The empirical outcome for a condition can then be estimated by\
\\[\
\\bar\{M\}_\{j\} = \\frac\{1\}\{Mn\} \\sum_\{i=1\}^\{M\}\\sum_\{r=1\}^\{n\} M(y_\{ijr\}).\
\\]\
where M is a evaluation metric.\
\\subsubsection*\{Step 6: Intermediate System Behavior\}\
Although the model is minimal, intermediate observations is to be logged.\\\\\
\\newline\
Record for each query:\
\\[\
(q_i,C_i,y_\{ijr\}).\
\\]\
\\newline\
If the system exposes citations, document identifiers, or other evidence references, these should also be recorded. This permits a distinction between\
retrieval and utilization of retrieved evidence. This distinction is important since a prompt may not alter which documents are retrieved but may alter how the generator uses those documents.\\\\\
\\section*\{Evaluation Metrics (To be discussed further; ChatGPT gave me some good intial ideas)\}\
\\subsubsection*\{1. Retrieval Utilization\}\
\
Let $C_i$ denote the retrieved context for query $q_i$, and let $E_i(y)\\subseteq C_i$ denote the subset of retrieved evidence explicitly or implicitly used to support the generated answer $y$.\\\\\
\\newline\
evidence recall measure is\
\\[\
ER(y) = \\frac\{|E(y)\\cap G|\}\{|G|\}.\
\\]\
evidence precision is measure is\
\\[\
EP(y) = \\frac\{|E(y)\\cap G|\}\{|E(y)|\}.\
\\]\
\
\\newline\
\
\
\\subsubsection*\{2. Hallucination Rate\}\
\
Let $H(y)$ denote the number of unsupported factual claims in output $y$, and let $N(y)$ denote the total number of factual claims.\\\\\
\\newline\
Hallucination rate defined as\
\\[\
\\boxed\{\
HR(y)=\\frac\{H(y)\}\{N(y)\}\
\}\
\\]\
where a claim is considered hallucinated when it cannot be supported by the available evidence.\\\\\
\\newline\
The aggregate hallucination rate for condition $S_j$ is\
\\[\
\\overline\{HR\}_j = \\frac\{1\}\{M\} \\sum_\{i=1\}^\{M\} HR(y_\{ij\}).\
\\]\
\\subsubsection*\{3. Answer Correctness\}\
Let\
\\[\
C_\{\\text\{ans\}\}(y,q)\\in[0,1]\
\\]\
denote the correctness of answer $y$ for query $q$.\\\\\
\\newline\
For binary evaluation,\
\\[\
C_\{\\text\{ans\}\}(y,q)\\in\\\{0,1\\\}.\
\\]\
\\newline\
The mean answer correctness under prompt $S_j$ is\
\\[\
\\overline\{C\}_j = \\frac\{1\}\{M\} \\sum_\{i=1\}^\{M\} C_\{\\text\{ans\}\}(y_\{ij\},q_i).\
\\]\
\
\\subsubsection*\{4. Evidence Faithfulness\}\
Let\
\\[\
C_\{\\text\{faith\}\}(y,C)\\in[0,1]\
\\]\
measure the extent to which factual claims in $y$ are supported by the retrieved context $C$.\\\\\
\\newline\
This metric is distinct from answer correctness. A response can be correct while failing to use the retrieved evidence, and it can be incorrect despite using evidence from the retrieved context.\\\\\
\\subsubsection*\{5. Constraint Adherence\}\
Let $K_j$ denote the set of behavioral constraints specified by prompt $S_j$.\\\\\
\\newline\
For each constraint $k\\in K_j$, define\
\\[\
A_k(y)\\in\\\{0,1\\\},\
\\]\
where $A_k(y)=1$ indicates that the output satisfies the constraint.\\\\\
\\newline\
The overall adherence score is\
\\[\
A(y,S_j) = \\frac\{1\}\{|K_j|\} \\sum_\{k\\in K_j\}A_k(y).\
\\]\
This directly measures whether a prompt intervention produces the intended behavioral change.\\\\\
\\subsubsection*\{6. Tone and Style Adherence\}\
If a prompt specifies an output style, define a style-adherence measure\
\\[\
T(y,S_j)\\in[0,1].\
\\]\
\\newline\
For instance, the prompt may specify that responses should be concise, formal, or structured. \\\\\
\
\\subsubsection*\{7. Abstention Behavior\}\
\
For questions where the available evidence is insufficient, define\
\\[\
A_\{\\text\{abs\}\}(y)\\in\\\{0,1\\\},\
\\]\
where $A_\{\\text\{abs\}\}(y)=1$ indicates an appropriate abstention.\\\\\
\\newline\
The abstention accuracy is\
\\[\
AA = \\frac\{\\text\{appropriate abstentions\}\}\{\\text\{insufficient-evidence queries\}\}.\
\\]\
\\newline\
This metric is particularly useful for evaluating prompts containing explicit uncertainty or abstention instructions.\
\\subsection*\{Output Length\}\
Log tokens for every generation. It confounds hallucination rate, utilization, and faithfulness simultaneously, and belongs in every model as a covariate.\
\\subsection*\{Self-consistency\}\
You run n repetitions but only average them. The dispersion across repetitions is an outcome in its own right: does a constraint reduce output variance? Pairwise semantic similarity among the n outputs, near-zero marginal cost.\
\\subsection*\{LLM as a Judge\}\
Correctness, faithfulness, hallucination, and style all require an LLM judge.\
\\subsection*\{Behavioral and Outcome Effects\}\
\
For a metric $M$, define the effect of prompt $S_j$ relative to the baseline $S_0$ as\
\\[\
\\boxed\{\
\\Delta M_j = \\overline\{M\}_j-\\overline\{M\}_0.\
\}\
\\]\
\\section*\{Drawing Conclusions\}\
\
The central analysis compares each experimental prompt condition with the baseline.\\\\\
\\newline\
For prompt $S_j$, the observed system can be represented as\
\\[\
F_j(q) = G(q,R(q;D);S_j).\
\\]\
\\newline\
The primary quantity is\
\\[\
\\Delta M_j = M(F_j)-M(F_0).\
\\]\
\\newline\
The interpretation of $\\Delta M_j$ depends on the metric.\
\\[\
\\Delta HR_j<0\
\\]\
would indicate a reduction in hallucination rate relative to the baseline.\
\\[\
\\Delta A_j>0\
\\]\
would indicate improved adherence to the constraints in the system prompt.\\\\\
\\newline\
The analysis should consider both statistical significance and practical significance. \\\\\
\\newline\
For paired observations, appropriate statistical tests can be selected according to the measurement scale and distribution of the data. For binary outcomes, paired proportion tests or logistic models may be appropriate; for continuous scores, paired tests or regression models may be used.\
\
\\section*\{Interpreting Prompt Effects\}\
\
The results should be interpreted according to three questions.\
\
\\subsubsection*\{Question 1: Does the prompt change behavior?\}\
This is evaluated primarily through variables such as instruction adherence, evidence utilization, tone/style adherence, and abstention behavior.\\\\\
\
\\subsubsection*\{Question 2: Does the behavioral change affect system outcomes?\}\
This is evaluated through variables such as answer correctness, evidence faithfulness, and hallucination rate.\\\\\
\\newline\
A prompt should not be considered as better only based on instruction adherence alone. It should also be evaluated according to the resulting behavior improves the quality and reliability of the system as a whole.\
\
\\subsubsection*\{Question 3: Do behavioral changes coincide with the intended behavior?\}\
The intermediate retrieval context $C$ provides an important control. Since $R$ is held fixed, differences in the final outputs can be examined as changes in the generator's use of the same evidence.\\\\\
\\newline\
The resulting mechanism is\
\\[\
\\boxed\{\
S \\rightarrow \\text\{behavior\} \\rightarrow \\text\{outcome\}.\
\}\
\\]\
\\newline\
The experiment identifies observable functional effects of the prompt within the controlled RAG architecture.\
\
\\newpage\
\\section*\{References\}\
\
\\begin\{itemize\}[leftmargin=*]\
\\item Cheng, Z., \\& Mastropaolo, A. (2026). \\textit\{An empirical study on the effects of system prompts in instruction-tuned models for code generation\}. arXiv preprint arXiv:2602.15228. \\url\{https://arxiv.org/abs/2602.15228\}\
\
\\item Mu, N., Lu, J., Lavery, M., \\& Wagner, D. A. (2025). \\textit\{A closer look at system prompt robustness\}. arXiv preprint arXiv:2502.12197. \\url\{https://arxiv.org/abs/2502.12197\}\
\
\\item Lu, S., Schuff, H., \\& Gurevych, I. (2024). How are prompts different in terms of sensitivity? In \\textit\{Proceedings of the 2024 Conference of the North American Chapter of the Association for Computational Linguistics: Human Language Technologies (NAACL-HLT)\}, 5833--5856. Association for Computational Linguistics. \\url\{https://aclanthology.org/2024.naacl-long.325/\}\
\
\\item Zhuo, J., Zhang, S., Fang, X., Duan, H., Lin, D., \\& Chen, K. (2024). ProSA: Assessing and understanding the prompt sensitivity of LLMs. In \\textit\{Findings of the Association for Computational Linguistics: EMNLP 2024\}. Association for Computational Linguistics. \\url\{https://aclanthology.org/2024.findings-emnlp.108/\}\
\
\\item Salinas, A., \\& Morstatter, F. (2024). The butterfly effect of altering prompts: How small changes and jailbreaks affect large language model performance. In \\textit\{Findings of the Association for Computational Linguistics: ACL 2024\}. Association for Computational Linguistics. \\url\{https://aclanthology.org/2024.findings-acl.275/\}\
\
\\item Zheng, M., Pei, J., Logeswaran, L., Lee, M., \\& Jurgens, D. (2024). When ``A helpful assistant'' is not really helpful: Personas in system prompts do not improve performances of large language models. In \\textit\{Findings of the Association for Computational Linguistics: EMNLP 2024\}, 15126--15154. Association for Computational Linguistics. \\url\{https://aclanthology.org/2024.findings-emnlp.888/\}\
\
\\item Ying, J., Cao, Y., Xiong, K., He, Y., Cui, L., \\& Liu, Y. (2024). Intuitive or dependent? Investigating LLMs' behavior style to conflicting prompts. In \\textit\{Proceedings of the 62nd Annual Meeting of the Association for Computational Linguistics (ACL 2024)\}, 4221--4246. Association for Computational Linguistics. \\url\{https://aclanthology.org/2024.acl-long.232/\}\
\
\\item Ruiz, F. J., \\& Cardona-Betancourt, V. (2026). Prompt carefully! ChatGPT displays rule-based insensitivity to contingencies. \\textit\{Journal of Contextual Behavioral Science, 41\}, 101012. \\url\{https://doi.org/10.1016/j.jcbs.2026.101012\}\
\
\\item Hua, A., Tang, K., Gu, C., Gu, J., Wong, E., \\& Qin, Y. (2025). Flaw or artifact? Rethinking prompt sensitivity in evaluating LLMs. In \\textit\{Proceedings of the 2025 Conference on Empirical Methods in Natural Language Processing\}. Association for Computational Linguistics. \\url\{https://aclanthology.org/2025.emnlp-main.1006/\}\
\
\\item Qin, R., Wang, Q., Wang, T., Wei, Z., \\& Shen, W. (2026). Evaluating and explaining prompt sensitivity of LLMs using interactions. arXiv preprint arXiv:2608.18539. \\url\{https://arxiv.org/abs/2608.18539\}\
\\end\{itemize\}\
\
\
\
\
\\end\{document\}\
}